\section{Fundamental Axioms of Stochastic Bidirectional Causality}

\subsection{Axiom I: Total Causal Structure}

Physical reality is represented by a total causal state,
\begin{equation}
\Psi_{\rm tot}(t),
\end{equation}
which contains both observable and non-observable causal degrees of freedom.

The observable universe is not assumed to exhaust the total causal structure.

\subsection{Axiom II: Finite Causal Accessibility}

Every physical state has access only to a finite causal sector of the total structure.

This finite accessibility is represented by a projection operator,
\begin{equation}
\mathcal P,
\end{equation}
such that
\begin{equation}
\Psi_{\rm obs}(t)
=
\mathcal P \Psi_{\rm tot}(t).
\end{equation}

The complementary sector is
\begin{equation}
\Psi_{\rm hid}(t)
=
(1-\mathcal P)\Psi_{\rm tot}(t).
\end{equation}

\subsection{Axiom III: Objective Observational Window}

The projection operator \(\mathcal P\) is not a subjective act of measurement.

It represents an objective causal window intrinsic to the structure of reality.

Thus,
\begin{equation}
\mathcal P^2=\mathcal P,
\end{equation}
and the decomposition
\begin{equation}
\Psi_{\rm tot}
=
\Psi_{\rm obs}
+
\Psi_{\rm hid}
\end{equation}
is physically meaningful.

\subsection{Axiom IV: Bidirectional Causal Coupling}

The observable and hidden sectors are not dynamically independent.

Their evolution is governed by coupled operators:
\begin{equation}
\frac{d\Psi_{\rm obs}}{dt}
=
\mathcal L_{oo}\Psi_{\rm obs}
+
\mathcal L_{oh}\Psi_{\rm hid},
\end{equation}
\begin{equation}
\frac{d\Psi_{\rm hid}}{dt}
=
\mathcal L_{ho}\Psi_{\rm obs}
+
\mathcal L_{hh}\Psi_{\rm hid}.
\end{equation}

The bidirectional structure is encoded in the nonzero couplings
\begin{equation}
\mathcal L_{oh}\neq 0,
\qquad
\mathcal L_{ho}\neq 0.
\end{equation}

\subsection{Axiom V: Emergent Memory}

Eliminating the hidden sector generates a nonlocal effective equation for the observable sector:
\begin{equation}
\frac{d\Psi_{\rm obs}}{dt}
=
\mathcal L_{\rm eff}\Psi_{\rm obs}
+
\int_{-\infty}^{t}
K(t,t')
\Psi_{\rm obs}(t')
\,dt'.
\end{equation}

Thus, memory is not fundamental.

It emerges from finite causal accessibility and hidden-sector elimination.

\subsection{Axiom VI: Stable Causal Screening}

The memory kernel must satisfy causality, finite normalization and stability:
\begin{equation}
K(t,t')=0
\qquad
{\rm for}
\qquad
t'>t,
\end{equation}
\begin{equation}
\int_{-\infty}^{t}
|K(t,t')|\,dt'
<
\infty,
\end{equation}
and
\begin{equation}
K(t,t')\rightarrow 0
\qquad
{\rm as}
\qquad
t-t'\rightarrow \infty.
\end{equation}

The minimal kernel satisfying these conditions is
\begin{equation}
K(t,t')
=
K_0 e^{-\lambda(t-t')}
\Theta(t-t').
\end{equation}

\subsection{Axiom VII: Infrared Effective Limit}

In the cosmological infrared limit, the memory response induces an effective perturbative deformation:
\begin{equation}
F_{\rm SBC}(z)
=
-\epsilon
\exp\left[-(k_0(1+z))^n\right].
\end{equation}

The minimal stable memory regime corresponds to
\begin{equation}
n\simeq 1.
\end{equation}

\subsection{Observational Constraint}

The empirical SBC/F3 sector selects
\begin{equation}
\epsilon\simeq 0.06,
\qquad
k_0\simeq 0.015-0.022,
\qquad
n\simeq 1,
\end{equation}
while preserving
\begin{equation}
H_{\rm SBC}(z)
\simeq
H_{\Lambda{\rm CDM}}(z).
\end{equation}

Thus, the observationally viable limit of SBC is
\begin{equation}
\Lambda{\rm CDM\ background}
+
{\rm SBC/F3\ perturbative\ infrared\ sector}.
\end{equation}
